\documentclass[12pt,a4paper]{report}

% ── Codificación y fuentes ──────────────────────────────────────────────────
\usepackage[utf8]{inputenc}
\usepackage[T1]{fontenc}
\usepackage[spanish]{babel}
\usepackage{lmodern}

% ── Márgenes y geometría ────────────────────────────────────────────────────
\usepackage[top=2.5cm, bottom=2.5cm, left=3cm, right=2.5cm]{geometry}

% ── Colores ─────────────────────────────────────────────────────────────────
\usepackage[dvipsnames,table]{xcolor}
\definecolor{navyblue}{RGB}{0,51,102}
\definecolor{lightgray}{RGB}{245,245,245}
\definecolor{warningyellow}{RGB}{255,193,7}
\definecolor{noteblue}{RGB}{13,110,253}
\definecolor{successgreen}{RGB}{25,135,84}

% ── Tablas ──────────────────────────────────────────────────────────────────
\usepackage{booktabs}
\usepackage{longtable}
\usepackage{array}
\usepackage{tabularx}
\usepackage{multirow}
\newcolumntype{L}[1]{>{\raggedright\arraybackslash}p{#1}}
\newcolumntype{C}[1]{>{\centering\arraybackslash}p{#1}}

% ── Gráficos y cajas ────────────────────────────────────────────────────────
\usepackage{graphicx}
\usepackage{float}
\usepackage{tcolorbox}
\tcbuselibrary{breakable,skins}

% ── Encabezados y pies de página ────────────────────────────────────────────
\usepackage{fancyhdr}
\pagestyle{fancy}
\fancyhf{}
\fancyhead[L]{\small\textcolor{navyblue}{\textbf{GEONAVAL} – Manual de Usuario}}
\fancyhead[R]{\small\textcolor{gray}{Sistema de Registro Fluvial y Monitoreo GPS}}
\fancyfoot[C]{\small\textcolor{gray}{Quibdó, Chocó – Colombia \quad|\quad Versión 1.0.0 – Mayo 2026}}
\fancyfoot[R]{\small Página \thepage}
\renewcommand{\headrulewidth}{0.4pt}
\renewcommand{\footrulewidth}{0.4pt}

% ── Hipervínculos ───────────────────────────────────────────────────────────
\usepackage[colorlinks=true, linkcolor=navyblue, urlcolor=noteblue]{hyperref}

% ── Títulos de secciones ────────────────────────────────────────────────────
\usepackage{titlesec}
\titleformat{\chapter}[hang]
  {\huge\bfseries\color{navyblue}}{\thechapter.}{0.5em}{}[\titlerule]
\titleformat{\section}[hang]
  {\large\bfseries\color{navyblue}}{\thesection}{0.5em}{}
\titleformat{\subsection}[hang]
  {\normalsize\bfseries\color{navyblue}}{\thesubsection}{0.5em}{}
\titlespacing{\chapter}{0pt}{10pt}{20pt}

% ── Listas ──────────────────────────────────────────────────────────────────
\usepackage{enumitem}
\setlist[itemize]{topsep=4pt,itemsep=2pt,leftmargin=1.5em}
\setlist[enumerate]{topsep=4pt,itemsep=2pt,leftmargin=1.5em}

% ── Cajas de Nota y Advertencia ─────────────────────────────────────────────
\newtcolorbox{notabox}{
  colback=noteblue!8, colframe=noteblue, fonttitle=\bfseries,
  title={\small Nota}, breakable, left=6pt, right=6pt, top=4pt, bottom=4pt
}
\newtcolorbox{advertenciabox}{
  colback=warningyellow!20, colframe=warningyellow!80!black, fonttitle=\bfseries,
  title={\small Advertencia}, breakable, left=6pt, right=6pt, top=4pt, bottom=4pt
}

% ── Tabla de contenido ──────────────────────────────────────────────────────
\setcounter{tocdepth}{2}

% ══════════════════════════════════════════════════════════════════════════════
\begin{document}

% ── Portada ─────────────────────────────────────────────────────────────────
\begin{titlepage}
  \centering
  \vspace*{2cm}
  {\Huge\bfseries\color{navyblue} GEONAVAL\par}
  \vspace{0.3cm}
  {\Large\color{gray} Control Fluvial\par}
  \vspace{1.2cm}
  \rule{0.8\textwidth}{1.5pt}
  \vspace{0.5cm}

  {\LARGE\bfseries MANUAL DE USUARIO\par}
  \vspace{0.4cm}
  {\large Sistema de Registro Fluvial y Monitoreo GPS\par}
  \vspace{0.2cm}
  {\normalsize Control y Seguridad del Transporte Fluvial\par}

  \vspace{1cm}
  \rule{0.8\textwidth}{0.8pt}
  \vspace{1cm}

  \begin{tabular}{@{}ll@{}}
    \toprule
    \textbf{Versión}   & 1.0.0 – Mayo 2026 \\
    \textbf{Materia}   & Ingeniería de Software \\
    \textbf{Ubicación} & Quibdó, Chocó – Colombia \\
    \textbf{Fecha}     & Mayo 2026 \\
    \bottomrule
  \end{tabular}

  \vfill
  {\small\color{gray} Quibdó, Chocó – Colombia \quad|\quad Versión 1.0.0 – Mayo 2026\par}
\end{titlepage}

% ── Tabla de contenido ──────────────────────────────────────────────────────
\tableofcontents
\newpage

% ══════════════════════════════════════════════════════════════════════════════
\chapter{Introducción}

\section{Descripción del Sistema}

GEONAVAL es un sistema web de Registro Fluvial y Monitoreo GPS diseñado para el control y la seguridad del transporte fluvial en la región de Quibdó, Chocó, Colombia. El sistema permite gestionar de forma integral todos los actores y procesos involucrados en la operación fluvial, incluyendo embarcaciones, propietarios, tripulación, pasajeros y viajes.

El sistema cuenta con monitoreo GPS en tiempo real, permitiendo a las autoridades y administradores hacer seguimiento a las embarcaciones activas sobre las rutas fluviales del río Atrato y sus afluentes.

\section{Propósito del Manual}

Este manual tiene como objetivo guiar a los usuarios del sistema GEONAVAL en el uso correcto de todas las funcionalidades disponibles. Está dirigido principalmente al perfil de Administrador, quien tiene acceso completo a todos los módulos del sistema.

\section{Alcance}

El presente manual cubre los siguientes módulos del sistema:

\begin{itemize}
  \item Inicio de sesión y autenticación
  \item Dashboard principal con indicadores clave
  \item Gestión de Embarcaciones
  \item Gestión de Propietarios
  \item Gestión de Tripulación
  \item Gestión de Pasajeros
  \item Viajes (Zarpe)
  \item Monitoreo GPS en Tiempo Real
  \item Compras y Tickets
  \item Reportes y Estadísticas
  \item Configuración del Sistema
\end{itemize}

\section{Convenciones del Manual}

\begin{table}[H]
  \centering
  \renewcommand{\arraystretch}{1.4}
  \begin{tabularx}{\textwidth}{@{}L{4.5cm} X@{}}
    \toprule
    \rowcolor{navyblue!15}
    \textbf{Elemento} & \textbf{Descripción} \\
    \midrule
    \textbf{Texto en negrita} & Indica nombres de botones, campos o elementos de interfaz \\
    \textbf{[Texto entre corchetes]} & Indica una acción que debe realizar el usuario \\
    \textbf{Nota:} & Información importante a tener en cuenta \\
    \textbf{Advertencia:} & Acción que puede tener consecuencias irreversibles \\
    \textbf{Flechas rojas (→)} & Señalan los elementos clave en cada captura de pantalla \\
    \bottomrule
  \end{tabularx}
  \caption{Convenciones tipográficas del manual}
\end{table}

% ══════════════════════════════════════════════════════════════════════════════
\chapter{Requisitos del Sistema}

\section{Requisitos Técnicos}

\begin{table}[H]
  \centering
  \renewcommand{\arraystretch}{1.4}
  \begin{tabularx}{\textwidth}{@{}L{4.5cm} X@{}}
    \toprule
    \rowcolor{navyblue!15}
    \textbf{Componente} & \textbf{Descripción} \\
    \midrule
    \textbf{Navegador web} & Google Chrome 90+, Mozilla Firefox 88+, Microsoft Edge 90+, Safari 14+ \\
    \textbf{Conexión a Internet} & Mínimo 2 Mbps (recomendado 5 Mbps o superior para GPS en tiempo real) \\
    \textbf{Resolución de pantalla} & Mínimo 1280 × 720 px (recomendado 1920 × 1080) \\
    \textbf{Sistema operativo} & Windows 10+, macOS 10.14+, Ubuntu 18.04+ \\
    \textbf{JavaScript} & Habilitado en el navegador \\
    \bottomrule
  \end{tabularx}
  \caption{Requisitos técnicos del sistema}
\end{table}

\section{Acceso al Sistema}

GEONAVAL es una aplicación web, por lo que no requiere instalación. El acceso se realiza directamente desde el navegador ingresando la URL proporcionada por el administrador de sistemas.

% ══════════════════════════════════════════════════════════════════════════════
\chapter{Inicio de Sesión}

\section{Acceso al Sistema}

Para ingresar a GEONAVAL, el usuario debe completar el formulario de autenticación que aparece al abrir el sistema. La siguiente imagen muestra la pantalla de inicio de sesión con sus elementos señalados:

\begin{figure}[H]
  \centering
  \fbox{\rule{0pt}{4cm}\hspace{10cm}}
  \caption{Pantalla de Inicio de Sesión de GEONAVAL}
  \label{fig:login}
\end{figure}

\section{Descripción de Elementos (ver Figura~\ref{fig:login})}

\begin{table}[H]
  \centering
  \renewcommand{\arraystretch}{1.4}
  \begin{tabularx}{\textwidth}{@{}L{4.5cm} X@{}}
    \toprule
    \rowcolor{navyblue!15}
    \textbf{Campo} & \textbf{Descripción} \\
    \midrule
    \textbf{① Correo Electrónico} & Campo donde se ingresa el correo registrado del usuario (ejemplo: \texttt{admin@geonaval.com}). \\
    \textbf{② Contraseña} & Campo de contraseña. El ícono de ojo permite mostrar u ocultar el texto. \\
    \textbf{③ Tipo de Usuario} & Menú desplegable para seleccionar el perfil: Administrador, Operador o Autoridad. \\
    \textbf{④ Iniciar Sesión} & Botón principal para autenticarse en el sistema. \\
    \textbf{⑤ Recuperar contraseña} & Enlace para restablecer la contraseña olvidada vía correo electrónico. \\
    \bottomrule
  \end{tabularx}
  \caption{Elementos de la pantalla de inicio de sesión}
\end{table}

\section{Pasos para Iniciar Sesión}

\begin{enumerate}
  \item Abra su navegador web preferido.
  \item Ingrese la URL del sistema GEONAVAL en la barra de direcciones.
  \item Complete el campo \textbf{Correo Electrónico} (①).
  \item Complete el campo \textbf{Contraseña} (②).
  \item Seleccione su \textbf{Tipo de Usuario} en el desplegable (③).
  \item Haga clic en \textbf{Iniciar Sesión} (④).
\end{enumerate}

\section{Tipos de Usuario}

\begin{table}[H]
  \centering
  \renewcommand{\arraystretch}{1.4}
  \begin{tabularx}{\textwidth}{@{}L{4cm} X@{}}
    \toprule
    \rowcolor{navyblue!15}
    \textbf{Perfil} & \textbf{Descripción} \\
    \midrule
    \textbf{Administrador} & Acceso completo a todos los módulos del sistema. \\
    \textbf{Operador} & Acceso a módulos operativos: compras, tickets, viajes y consulta de embarcaciones. \\
    \textbf{Autoridad} & Acceso de solo lectura a monitoreo, reportes y alertas. \\
    \bottomrule
  \end{tabularx}
  \caption{Perfiles de usuario del sistema}
\end{table}

% ══════════════════════════════════════════════════════════════════════════════
\chapter{Dashboard Principal}

\section{Descripción General}

Una vez autenticado, el sistema redirige al Dashboard Principal. La siguiente imagen muestra sus secciones principales:

\begin{figure}[H]
  \centering
  \fbox{\rule{0pt}{4cm}\hspace{10cm}}
  \caption{Dashboard Principal de GEONAVAL}
  \label{fig:dashboard}
\end{figure}

\section{Descripción de Elementos (ver Figura~\ref{fig:dashboard})}

\begin{table}[H]
  \centering
  \renewcommand{\arraystretch}{1.4}
  \begin{tabularx}{\textwidth}{@{}L{4.5cm} X@{}}
    \toprule
    \rowcolor{navyblue!15}
    \textbf{Elemento} & \textbf{Descripción} \\
    \midrule
    \textbf{① Menú de navegación} & Panel lateral izquierdo con acceso a todos los módulos del sistema. \\
    \textbf{② Embarcaciones Activas} & KPI con el número de embarcaciones operativas y variación mensual. \\
    \textbf{③ Viajes en Curso} & KPI con el número de viajes activos y variación mensual. \\
    \textbf{④ Perfil de usuario} & Área superior derecha con nombre, rol y opciones de sesión. \\
    \textbf{⑤ Mapa GPS en tiempo real} & Vista del mapa fluvial con posición de embarcaciones activas. \\
    \textbf{⑥ Próxima Llegada} & Panel con hora estimada y datos de la siguiente embarcación en llegar. \\
    \bottomrule
  \end{tabularx}
  \caption{Elementos del Dashboard Principal}
\end{table}

\section{Indicadores Clave (KPIs)}

Los cuatro KPIs de la parte superior muestran: \textbf{Embarcaciones Activas}, \textbf{Viajes en Curso}, \textbf{Pasajeros Transportados} y \textbf{Alertas Activas}, todos con comparativo frente al mes anterior.

% ══════════════════════════════════════════════════════════════════════════════
\chapter{Gestión de Embarcaciones}

\section{Descripción}

El módulo de Gestión de Embarcaciones permite registrar, consultar y administrar todas las embarcaciones fluviales. La siguiente imagen muestra la vista principal del módulo:

\begin{figure}[H]
  \centering
  \fbox{\rule{0pt}{4cm}\hspace{10cm}}
  \caption{Módulo de Gestión de Embarcaciones}
  \label{fig:embarcaciones}
\end{figure}

\section{Descripción de Elementos (ver Figura~\ref{fig:embarcaciones})}

\begin{table}[H]
  \centering
  \renewcommand{\arraystretch}{1.4}
  \begin{tabularx}{\textwidth}{@{}L{5cm} X@{}}
    \toprule
    \rowcolor{navyblue!15}
    \textbf{Elemento} & \textbf{Descripción} \\
    \midrule
    \textbf{① Nueva Embarcación} & Botón para abrir el formulario de registro de una nueva embarcación. \\
    \textbf{② KPIs del módulo} & Resumen de totales: embarcaciones registradas, operativas, en mantenimiento y viajes. \\
    \textbf{③ Estado: Operativa} & Etiqueta verde que indica que la embarcación está disponible para viajes. \\
    \textbf{④ Estado: En Mantenimiento} & Etiqueta naranja que indica que la embarcación está temporalmente fuera de servicio. \\
    \textbf{⑤ Ver detalles} & Enlace para consultar toda la información de una embarcación específica. \\
    \bottomrule
  \end{tabularx}
  \caption{Elementos del módulo de Gestión de Embarcaciones}
\end{table}

\section{Registrar Nueva Embarcación}

\begin{enumerate}
  \item Haga clic en \textbf{+ Nueva Embarcación} (①).
  \item Complete el formulario: nombre, tipo, propietario y estado inicial.
  \item Haga clic en \textbf{Guardar} para registrar.
\end{enumerate}

% ══════════════════════════════════════════════════════════════════════════════
\chapter{Gestión de Propietarios}

\section{Descripción}

Este módulo permite administrar el registro de propietarios de embarcaciones, tanto personas naturales como empresas. La imagen a continuación muestra la vista principal:

\begin{figure}[H]
  \centering
  \fbox{\rule{0pt}{4cm}\hspace{10cm}}
  \caption{Módulo de Gestión de Propietarios}
  \label{fig:propietarios}
\end{figure}

\section{Descripción de Elementos (ver Figura~\ref{fig:propietarios})}

\begin{table}[H]
  \centering
  \renewcommand{\arraystretch}{1.4}
  \begin{tabularx}{\textwidth}{@{}L{5cm} X@{}}
    \toprule
    \rowcolor{navyblue!15}
    \textbf{Elemento} & \textbf{Descripción} \\
    \midrule
    \textbf{① Nuevo Propietario} & Botón para registrar un nuevo propietario en el sistema. \\
    \textbf{② Tipo: Natural / Empresa} & Clasificación del propietario como persona física o jurídica. \\
    \textbf{③ Ver / Editar / Eliminar} & Íconos de acción disponibles por cada registro. \\
    \textbf{④ Embarcaciones asociadas} & Contador que muestra cuántas embarcaciones tiene cada propietario. \\
    \bottomrule
  \end{tabularx}
  \caption{Elementos del módulo de Gestión de Propietarios}
\end{table}

\begin{advertenciabox}
Al eliminar un propietario, esta acción es \textbf{irreversible}. Verifique que no tenga embarcaciones activas antes de proceder.
\end{advertenciabox}

% ══════════════════════════════════════════════════════════════════════════════
\chapter{Gestión de Tripulación}

\section{Descripción}

El módulo de Tripulación permite administrar el personal operativo de las embarcaciones. La siguiente imagen muestra la vista principal:

\begin{figure}[H]
  \centering
  \fbox{\rule{0pt}{4cm}\hspace{10cm}}
  \caption{Módulo de Gestión de Tripulación}
  \label{fig:tripulacion}
\end{figure}

\section{Descripción de Elementos (ver Figura~\ref{fig:tripulacion})}

\begin{table}[H]
  \centering
  \renewcommand{\arraystretch}{1.4}
  \begin{tabularx}{\textwidth}{@{}L{5cm} X@{}}
    \toprule
    \rowcolor{navyblue!15}
    \textbf{Elemento} & \textbf{Descripción} \\
    \midrule
    \textbf{① Nuevo Tripulante} & Botón para registrar un nuevo miembro de la tripulación. \\
    \textbf{② KPIs del módulo} & Total de tripulantes, operadores fluviales, personal auxiliar y activos. \\
    \textbf{③ Rol del tripulante} & Cargo asignado: Operador Fluvial, Auxiliar de Pasajeros, Motorista, etc. \\
    \textbf{④ Estado: Activo} & Etiqueta verde que indica disponibilidad del tripulante. \\
    \textbf{⑤ Acciones} & Íconos para Ver, Editar y Eliminar cada registro de tripulante. \\
    \bottomrule
  \end{tabularx}
  \caption{Elementos del módulo de Gestión de Tripulación}
\end{table}

\section{Registrar Nuevo Tripulante}

\begin{enumerate}
  \item Haga clic en \textbf{+ Nuevo Tripulante} (①).
  \item Complete nombre completo, rol, embarcación asignada y horario.
  \item Defina el estado inicial (Activo/Inactivo) y haga clic en \textbf{Guardar}.
\end{enumerate}

% ══════════════════════════════════════════════════════════════════════════════
\chapter{Gestión de Compras y Tickets}

\section{Descripción}

Este módulo centraliza la venta de pasajes y la gestión de tickets de transporte fluvial. La imagen siguiente muestra la vista principal:

\begin{figure}[H]
  \centering
  \fbox{\rule{0pt}{4cm}\hspace{10cm}}
  \caption{Módulo de Compras y Tickets}
  \label{fig:compras}
\end{figure}

\section{Descripción de Elementos (ver Figura~\ref{fig:compras})}

\begin{table}[H]
  \centering
  \renewcommand{\arraystretch}{1.4}
  \begin{tabularx}{\textwidth}{@{}L{5cm} X@{}}
    \toprule
    \rowcolor{navyblue!15}
    \textbf{Elemento} & \textbf{Descripción} \\
    \midrule
    \textbf{① Nueva Compra} & Botón para registrar una nueva venta de ticket. \\
    \textbf{② KPIs del módulo} & Ventas del día, total recaudado, tickets confirmados y pendientes. \\
    \textbf{③ Barra de búsqueda} & Permite buscar tickets por número, pasajero o documento. \\
    \textbf{④ Filtrar por Fecha} & Abre un selector de rango de fechas para filtrar los registros. \\
    \textbf{⑤ Exportar} & Descarga el listado de tickets en formato Excel. \\
    \textbf{⑥ Estado Confirmado} & Etiqueta azul que indica que el pago ha sido verificado. \\
    \textbf{⑦ Estado Pendiente} & Etiqueta amarilla que indica que el ticket aún requiere confirmación. \\
    \bottomrule
  \end{tabularx}
  \caption{Elementos del módulo de Compras y Tickets}
\end{table}

\section{Crear Nueva Compra}

\begin{enumerate}
  \item Haga clic en \textbf{+ Nueva Compra} (①).
  \item Seleccione el pasajero o regístrelo si es nuevo.
  \item Elija la ruta, fecha, asiento y método de pago.
  \item Confirme la transacción para generar el ticket.
\end{enumerate}

% ══════════════════════════════════════════════════════════════════════════════
\chapter{Monitoreo GPS en Tiempo Real}

\section{Descripción}

El módulo de Monitoreo GPS permite hacer seguimiento en tiempo real de las embarcaciones activas. La imagen siguiente muestra la vista del mapa:

\begin{figure}[H]
  \centering
  \fbox{\rule{0pt}{4cm}\hspace{10cm}}
  \caption{Monitoreo GPS en Tiempo Real}
  \label{fig:gps}
\end{figure}

\section{Descripción de Elementos (ver Figura~\ref{fig:gps})}

\begin{table}[H]
  \centering
  \renewcommand{\arraystretch}{1.4}
  \begin{tabularx}{\textwidth}{@{}L{5.5cm} X@{}}
    \toprule
    \rowcolor{navyblue!15}
    \textbf{Elemento} & \textbf{Descripción} \\
    \midrule
    \textbf{① Indicador EN VIVO} & Punto rojo parpadeante que confirma que los datos son en tiempo real. \\
    \textbf{② Embarcación en curso (verde)} & Punto verde que representa una embarcación navegando actualmente. \\
    \textbf{③ Embarcación programada (azul)} & Punto azul que indica un viaje confirmado pendiente de inicio. \\
    \textbf{④ Leyenda del mapa} & Guía visual de colores: verde = en curso, azul = programado, naranja = alerta. \\
    \bottomrule
  \end{tabularx}
  \caption{Elementos del módulo de Monitoreo GPS}
\end{table}

\begin{notabox}
Si una embarcación no aparece en el mapa, es posible que su dispositivo GPS esté apagado o sin señal. Verifique el dispositivo a bordo de la embarcación.
\end{notabox}

% ══════════════════════════════════════════════════════════════════════════════
\chapter{Reportes y Estadísticas}

\section{Filtros y Gráficas}

El módulo de Reportes ofrece herramientas de análisis estadístico. La imagen siguiente muestra los filtros y las gráficas disponibles:

\begin{figure}[H]
  \centering
  \fbox{\rule{0pt}{4cm}\hspace{10cm}}
  \caption{Reportes: Filtros de Búsqueda y Gráficas Estadísticas}
  \label{fig:reportes1}
\end{figure}

\section{Descripción de Elementos -- Figura~\ref{fig:reportes1}}

\begin{table}[H]
  \centering
  \renewcommand{\arraystretch}{1.4}
  \begin{tabularx}{\textwidth}{@{}L{5cm} X@{}}
    \toprule
    \rowcolor{navyblue!15}
    \textbf{Elemento} & \textbf{Descripción} \\
    \midrule
    \textbf{① Filtros de búsqueda} & Permite definir Fecha Inicio, Fecha Fin, Embarcación, Ruta y Estado para acotar el reporte. \\
    \textbf{② Aplicar / Limpiar} & Botones para aplicar los filtros seleccionados o restablecer los valores por defecto. \\
    \textbf{③ Viajes/Pasajeros} & Gráfica de líneas con evolución mensual de viajes y pasajeros transportados. \\
    \textbf{④ Uso de Embarcaciones} & Gráfica de barras con viajes y horas de operación por embarcación. \\
    \textbf{⑤ Tooltip interactivo} & Al pasar el cursor sobre la gráfica se muestran los valores exactos del mes. \\
    \bottomrule
  \end{tabularx}
  \caption{Elementos de filtros y gráficas del módulo de Reportes}
\end{table}

\section{Descarga de Reportes}

La imagen siguiente muestra los tipos de reportes disponibles y las opciones de descarga:

\begin{figure}[H]
  \centering
  \fbox{\rule{0pt}{4cm}\hspace{10cm}}
  \caption{Reportes: Tipos de Reportes y Descarga}
  \label{fig:reportes2}
\end{figure}

\section{Descripción de Elementos -- Figura~\ref{fig:reportes2}}

\begin{table}[H]
  \centering
  \renewcommand{\arraystretch}{1.4}
  \begin{tabularx}{\textwidth}{@{}L{5cm} X@{}}
    \toprule
    \rowcolor{navyblue!15}
    \textbf{Elemento} & \textbf{Descripción} \\
    \midrule
    \textbf{① Tipo de reporte} & Tarjeta con el nombre y descripción del reporte (Viajes del Mes, Pasajeros, Rutas, etc.). \\
    \textbf{② Descargar PDF} & Descarga el reporte en formato PDF (solo lectura, ideal para compartir). \\
    \textbf{③ Descargar Excel} & Descarga el reporte en formato Excel (editable, ideal para análisis). \\
    \textbf{④ Resumen estadístico} & Indicadores consolidados del período: total viajes, pasajeros, horas, tasa de finalización. \\
    \textbf{⑤ Pasajeros totales} & Cifra total de pasajeros transportados en el período seleccionado. \\
    \bottomrule
  \end{tabularx}
  \caption{Elementos de descarga del módulo de Reportes}
\end{table}

% ══════════════════════════════════════════════════════════════════════════════
\chapter{Configuración del Sistema}

\section{Descripción}

El módulo de Configuración permite gestionar el perfil, notificaciones y seguridad del sistema. La imagen siguiente muestra su vista principal:

\begin{figure}[H]
  \centering
  \fbox{\rule{0pt}{4cm}\hspace{10cm}}
  \caption{Módulo de Configuración del Sistema}
  \label{fig:config}
\end{figure}

\section{Descripción de Elementos (ver Figura~\ref{fig:config})}

\begin{table}[H]
  \centering
  \renewcommand{\arraystretch}{1.4}
  \begin{tabularx}{\textwidth}{@{}L{5.5cm} X@{}}
    \toprule
    \rowcolor{navyblue!15}
    \textbf{Elemento} & \textbf{Descripción} \\
    \midrule
    \textbf{① Editar Perfil} & Botón para modificar los datos del usuario activo (nombre, correo, etc.). \\
    \textbf{② Datos del usuario activo} & Muestra nombre, correo, rol, último acceso y fecha de registro. \\
    \textbf{③ Notificaciones activas} & Casillas marcadas para alertas de viajes, GPS y embarcaciones. \\
    \textbf{④ Notificación desactivada} & Casilla desmarcada para reportes automáticos (desactivados por defecto). \\
    \textbf{⑤ Opciones de Seguridad} & Accesos a: cambiar contraseña, autenticación 2FA, dispositivos y historial de sesiones. \\
    \bottomrule
  \end{tabularx}
  \caption{Elementos del módulo de Configuración del Sistema}
\end{table}

% ══════════════════════════════════════════════════════════════════════════════
\chapter{Glosario de Términos}

\begin{table}[H]
  \centering
  \renewcommand{\arraystretch}{1.5}
  \begin{tabularx}{\textwidth}{@{}L{4.5cm} X@{}}
    \toprule
    \rowcolor{navyblue!15}
    \textbf{Término} & \textbf{Descripción} \\
    \midrule
    \textbf{Dashboard} & Pantalla principal del sistema con indicadores y resúmenes visuales de la operación. \\
    \textbf{Embarcación} & Cualquier nave fluvial registrada: ferry, lancha, bote, etc. \\
    \textbf{GPS} & Sistema de Posicionamiento Global. Permite rastrear la ubicación de las embarcaciones en tiempo real. \\
    \textbf{KPI} & Indicador Clave de Rendimiento. Métrica para evaluar el desempeño del sistema. \\
    \textbf{Monitoreo en tiempo real} & Seguimiento continuo y actualizado de posición y estado de embarcaciones. \\
    \textbf{Propietario} & Persona natural o jurídica dueña de una o más embarcaciones registradas. \\
    \textbf{Ticket} & Documento digital que representa el derecho de un pasajero a viajar en una embarcación. \\
    \textbf{Tripulación} & Personal operativo de una embarcación: operadores, auxiliares, motoristas. \\
    \textbf{Zarpe} & Autorización formal para que una embarcación parta hacia su destino. \\
    \textbf{Ruta fluvial} & Trayecto definido por el río que une dos puntos (origen y destino). \\
    \textbf{NIT} & Número de Identificación Tributaria. Código fiscal de empresas en Colombia. \\
    \textbf{Alerta} & Notificación ante una situación anormal: desviación de ruta, fallo GPS, etc. \\
    \bottomrule
  \end{tabularx}
  \caption{Glosario de términos del sistema GEONAVAL}
\end{table}

% ══════════════════════════════════════════════════════════════════════════════
\chapter{Preguntas Frecuentes (FAQ)}

\subsection*{¿Qué hago si olvidé mi contraseña?}
En la pantalla de inicio de sesión, haga clic en el enlace \textit{«¿Olvidaste tu contraseña?»}. El sistema enviará un correo con instrucciones para restablecerla.

\subsection*{¿Por qué no aparece una embarcación en el mapa GPS?}
Si una embarcación no es visible en el mapa, puede deberse a que su dispositivo GPS está apagado, sin señal, o fuera del rango de cobertura. Verifique el dispositivo a bordo.

\subsection*{¿Cómo exporto un reporte en Excel?}
En el módulo de Reportes, seleccione el tipo de reporte deseado y haga clic en el botón \textbf{Excel} (verde). El archivo se descargará automáticamente.

\subsection*{¿Puedo eliminar un propietario que tiene embarcaciones asociadas?}
No se recomienda. Antes de eliminar un propietario, reasigne o elimine sus embarcaciones. El sistema podría impedir la eliminación si detecta registros relacionados activos.

\subsection*{¿Cómo sé si un ticket está confirmado?}
En el módulo de Compras y Tickets, la columna \textbf{Estado} muestra \textit{Confirmado} en azul o \textit{Pendiente} en amarillo para identificar el estado de cada ticket fácilmente.

\subsection*{¿Quién puede acceder al módulo de Configuración?}
Solo los usuarios con perfil \textbf{Administrador} tienen acceso completo. Los perfiles Operador y Autoridad tienen acceso restringido.

% ══════════════════════════════════════════════════════════════════════════════
\chapter{Contacto y Soporte Técnico}

Para soporte técnico, reportes de errores o consultas sobre el sistema GEONAVAL, comuníquese a través de los siguientes canales:

\begin{table}[H]
  \centering
  \renewcommand{\arraystretch}{1.5}
  \begin{tabularx}{\textwidth}{@{}L{4.5cm} X@{}}
    \toprule
    \rowcolor{navyblue!15}
    \textbf{Canal} & \textbf{Detalle} \\
    \midrule
    \textbf{Teléfono} & +57 (4) 670-1234 \\
    \textbf{Correo electrónico} & \texttt{info@geonaval.gov.co} \\
    \textbf{Dirección} & Calle Principal \#10-20, Quibdó, Chocó, Colombia \\
    \textbf{Horario de atención} & Lunes a Viernes: 8:00 AM – 6:00 PM \quad|\quad Sábado: 8:00 AM – 12:00 PM \\
    \textbf{Emergencias GPS} & Disponible 24/7 para incidentes críticos de monitoreo \\
    \bottomrule
  \end{tabularx}
  \caption{Canales de contacto y soporte técnico}
\end{table}

\vspace{1cm}
\begin{tcolorbox}[colback=navyblue!8, colframe=navyblue, title={\textbf{Información Legal}}]
GEONAVAL es un Sistema Oficial autorizado por la Capitanía del Puerto.\\
Versión 1.0.0 – Mayo 2026. Todos los derechos reservados.
\end{tcolorbox}

\vfill
\begin{center}
  \small\color{gray} Quibdó, Chocó – Colombia \quad|\quad Versión 1.0.0 – Mayo 2026
\end{center}

\end{document}
